\irishDanceName{The Walls of Limerick - Burchenal 1929}{Reel}{Progresivní}{Longway of sets}{4+2n}
% Name, tune, type, formation, number of dancers

{\large Progresivní reel pro 4 + 2n tanečníků -- longway vždy dvou párů naproti sobě}\hfill  

Dva páry stojí proti sobě, ženy napravo od partnerů.
Tanec je progresivní, takže každý pár postupuje stále stejným směrem, dokud nedorazí na konec longwaye.
Jako melodie je doporučena \uv{Limerick Lasses}, ale jakýkoliv AABB reel (2x8+2x8 taktů) by měl stačit.

Toto je \uv{archivní} verze tance tak, jak je uvedena ve sbírce \Source{Burchenal1929}. 

\HRule
\HRule
\vfill
\begin{description}
    \danceFigure{A}{Advance \& Retire \bars{4}}{
        \item[1–4]Advance \& Retire, neopakovat
    }

    \danceFigure{B}{Výměny dámy, pak pánové \bars{4}}{
        \item[1–2]Dámy výměna míst sidestepem vpravo (míjet čelem), bez záskoků
        \item[3–4]Pánové výměna míst sidestepem vlevo (míjet čelem), bez záskoků
    }

    \danceFigure{C}{Dance with Opposite \bars{8}}{
        \item[1–2]Pán s protější dámou za obě ruce v kříženém držení sidestep \uv{ven} (do pánova leva)
        \item[3-4]Zátočka o půlku na dvě promenády
        \item[5–6]Sidestep zpět
        \item[7-8]Zátočka o půlku do výchozích míst této figury 
    }

    \danceFigure{D}{Dance Around \bars{8}}{
        \item[1–8]Swing v kříženém držení o celé kolo, na konci koukat na další pár (zády k~předchozímu)
    }
\end{description}
